\documentclass[12pt]{article}
\usepackage[utf8]{inputenc}
\usepackage[T1]{fontenc}
\usepackage{hyperref}
\title{Miwok, Plains Language Resources for K-12 Curriculum Development}
\author{Compiled by Claude Sonnet 5.5 (claude-sonnet-5-5) for Andr\'as Kornai}
\date{}
\begin{document}
\maketitle
\section{Introduction}
Plains Miwok (ISO 639-3 pmw) was spoken east of the Sacramento River, from near Sacramento south to Rio Vista and into the Sierra foothills. It has no first-language speakers today; tribal members and activists pursue revitalization. The documentation rests chiefly on Catherine Callaghan's work.
\section{Primary Educational Resources}
No school materials were found in search results.
\section{Linguistic and Reference Materials}
\begin{itemize}
\item Callaghan, C. A. (1984), \emph{Plains Miwok Dictionary}, UC Publications in Linguistics 105 (309 pp.; English-Miwok and Miwok-English): record at \url{https://wiki.ubc.ca/Documentation:RelLex/Plains_Miwok_Dictionary}
\item University of Pennsylvania library catalogue record: \url{https://find.library.upenn.edu/catalog/992634733503681}
\item Wikipedia article: \url{https://en.wikipedia.org/wiki/Plains_Miwok_language}
\item Berkeley Linguistics language page: \url{https://cla.berkeley.edu/languages/plains-miwok.php}
\end{itemize}
\section{Religious and Community Resources}
None identified in search results.
\section{Digital and Audio Resources}
\begin{itemize}
\item Finding aid to the Catherine A. Callaghan papers (Bancroft Library, including Plains Miwok field notes): \url{https://oac.cdlib.org/findaid/ark:/13030/c8794bk2}
\item American Philosophical Society guide entry, Miwok materials, Harvey Pitkin Papers: \url{https://indigenousguide.amphilsoc.org/entry/38846}
\end{itemize}
\section{Curriculum Development Context}
A heritage curriculum would have to be rebuilt from the dictionary and archival notes, in partnership with Miwok tribal members and language revitalization activists. No existing school programme was found.
\end{document}